% H9-T14: constant finite-field varieties after transcendental base extension.
% Run JAC1964-20261121; /root/h9_height; GPT-6 Astra Pro.
% Input P15 uses the Tate consequence032 Cor8.3 and Kahn2003 Th1.10.

\subsection*{H9-T14. Constant finite-field varieties and the full primitive height domain}

\begin{lemma}[Spreading algebraic equivalence]\label{h9:spread}
Let $k_0$ be an algebraically closed field and let $X_0/k_0$ be
smooth projective. Fix a codimension $r$ and a prime $\ell$ invertible
in $k_0$, and use rational $\ell$-adic cycle classes. These classes are
compatible with smooth proper base change and are unchanged under
extension of algebraically closed fields.
Suppose
\[
 \operatorname{CH}^r(X_0)_{\rm hom,\mathbf Q}=0.
                                                               \tag{H9.48}
\]
Then, for every extension field $F/k_0$, including imperfect fields,
\[
 \operatorname{CH}^r((X_0)_F)_{\rm hom,\mathbf Q}
       =\operatorname{CH}^r((X_0)_F)_{\rm alg,\mathbf Q}.       \tag{H9.49}
\]
Homological triviality on the left is geometric. Algebraic triviality
on the right means algebraic equivalence already over $F$, with
rational coefficients, not merely after passage to $\overline F$.
\end{lemma}

\begin{proof}
Only the inclusion from left to right requires proof. Let $z$ be a
rational cycle class on $(X_0)_F$ whose geometric class is zero.
Choose a cycle representative and a finitely generated subfield
$F_0/k_0$ of $F$ carrying that cycle. Since $k_0$ is perfect,
$F_0/k_0$ is separably generated. It has a smooth geometrically
integral affine model $S/k_0$ with $k_0(S)=F_0$.

Taking closures of the finitely many cycle components and shrinking
$S$ if necessary gives a relative cycle
\[
 Z\in\operatorname{CH}^r(X_0\times S)_{\mathbf Q}
\]
whose generic fibre pulls back to $z$, and whose components are flat
over $S$. All fibres are interpreted by refined Gysin pullback; on
this flat locus these are the usual cycle-theoretic fibres. The
global cycle class of $Z$ has an image under the edge restriction map
\[
 H^{2r}(X_0\times S,\mathbf Q_\ell(r))
  \longrightarrow
 H^0\!\left(S,R^{2r}\operatorname{pr}_{S*}\mathbf Q_\ell(r)\right).
                                                               \tag{H9.50}
\]
Here the displayed sheaf is constant, with fibre
$H^{2r}(X_0,\mathbf Q_\ell(r))$, because the family is the product
family. The image of the class at a geometric generic point is zero:
its pullback to $\overline F$ is the zero class of $z$, and extension
between algebraically closed fields preserves this cohomology.
Since $S$ is connected, the section in (H9.50) is zero everywhere.
Compatibility of cycle classes with Gysin restriction therefore
gives $\operatorname{cl}(Z_s)=0$ for every geometric fibre.

Choose a closed point $s\in S(k_0)$; such a point exists because
$k_0$ is algebraically closed and $S$ is a nonempty scheme of finite
type. Hypothesis (H9.48) gives $Z_s=0$ in
$\operatorname{CH}^r(X_0)_{\mathbf Q}$.

There are now two \emph{$F$-rational} points of the smooth
geometrically integral parameter space $S_F$: the point $s_F$ and
the point $t$ induced by $F_0=k_0(S)\hookrightarrow F$. Although
$t$ maps to the generic point of $S$, it is a rational closed point
of $S_F$. By construction,
\[
 (Z_F)_t=z,\qquad (Z_F)_{s_F}=(Z_s)_F=0.
                                                               \tag{H9.51}
\]
Thus (H9.51) is already a smooth connected parameter-space witness
for algebraic triviality over $F$.

To verify the curve formulation explicitly, if $\dim S>0$, choose
a smooth geometrically integral curve in the quasi-projective smooth
variety $S_F$ through $t$ and $s_F$. The required Bertini statement
is Achter--Casalaina-Martin--Vial, \emph{Parameter spaces for
algebraic equivalence}, Theorem A.1(1),(4), author manuscript p.~17.
Pulling $Z_F$ back to that curve preserves (H9.51). If $F$ is
algebraically closed, normalize its projective compactification and
close the cycle family there; the resulting smooth projective
parameter curve has the same two fibres. Over an imperfect $F$, a
smooth quasi-projective parameter curve suffices for algebraic
equivalence. If a smooth projective curve is desired, Lemma 3.8(2)
of the same source, p.~6, supplies it after multiplication by a power
of the characteristic, harmless with rational coefficients. If
$\dim S=0$, then $F_0=k_0$ and (H9.48) already says $z=0$.
This proves (H9.49).
\end{proof}

\paragraph{The numerical variant, without a numerical/homological conjecture.}
Suppose instead that
$\operatorname{CH}^r(X_0)_{\rm num,\mathbf Q}=0$.
Then for every extension $F/k_0$ one has
\[
 \operatorname{CH}^r((X_0)_F)_{\rm num,\mathbf Q}
   =\operatorname{CH}^r((X_0)_F)_{\rm hom,\mathbf Q}
   =\operatorname{CH}^r((X_0)_F)_{\rm alg,\mathbf Q}.         \tag{H9.49a}
\]
Here numerical triviality may even be tested only against cycles
defined over $F$; the proof shows it implies geometric numerical
triviality as well. To see this, spread a numerically trivial $z$
as above. For \emph{each} complementary cycle $w$ on $X_0$, form
\[
 \operatorname{pr}_{S*}
       \bigl(Z\cdot\operatorname{pr}_{X_0}^*w\bigr)
       =d_w[S]\quad\text{in }\operatorname{CH}^0(S)_{\mathbf Q}.
\]
The intersection product exists because $X_0\times S$ is smooth;
no proper-intersection assumption on the chosen representatives is
needed. The projection is proper, and $S$ is integral, so a
codimension-zero class has precisely this form. Compatibility of
proper pushforward with refined Gysin restriction says that every
fibre degree $\deg(Z_s\cdot w)$ equals $d_w$. The same assertion
at the generic point is flat base change. Pulling that generic
intersection to $F$ preserves its degree, and numerical triviality
of $z$ gives $d_w=0$.

Thus the closed $k_0$-fibre $Z_s$ is numerically trivial, since the
argument applies to every complementary $w$ on $X_0$. It is zero
in rational Chow by the new hypothesis. Equation (H9.51) again
proves that $z$ is algebraically trivial over $F$. Finally,
algebraic triviality implies geometric homological triviality,
which implies numerical triviality. These inclusions prove all
equalities in (H9.49a). This argument uses intersection degrees,
not a standard conjecture identifying numerical and homological
equivalence.

\begin{theorem}[H9-T14]\label{h9:T14}
Put $k_0=\overline{\mathbf F}_p$. Accept the Tate consequence of
032, Corollary 8.3, and the established Kahn input in JAC-P15.

\begin{enumerate}
\item
Let $X_0/k_0$ be smooth projective with geometric rational Chow
motive of abelian type. For every extension field $F/k_0$ and every
codimension $r$ one has the three equalities (H9.49a). In particular
this applies to every
abelian variety over $k_0$ and to every product of curves over $k_0$.
For an abelian variety over $F$ geometrically isogenous to
$(A_0)_{\overline F}$, with $A_0/k_0$ abelian, the same equalities
hold. Consequently an integral
cycle on any of these varieties is numerically trivial if and only
if a nonzero multiple is algebraically trivial: this is customary
$\tau$-triviality, including after transcendental extension.

\item
Let $k/k_0$ be algebraically closed, let $C/k$ be a smooth projective
connected curve, and put $K=k(C)$. Let $A_0/k_0$ be an abelian
variety of dimension $2m-1$, and write $A=(A_0)_k$. Choose an ample
class on $A$ and let
\[
 B=P^m(A)_{\rm prim},\qquad
 \ell_B=\frac{(-1)^m}{2}\Delta^*d_B.
\]
Accept also 032, Corollary 8.4. Then $\ell_B$ is a rational
polarization by JAC-T4. Let
\[
 a_K:\operatorname{CH}^m(A_K)_{\rm hom,\mathbf Q}
          \longrightarrow P^m(A)_K(K)_{\mathbf Q}
\]
be the incidence map, which is now defined on the entire displayed
homological group by part (1). Its primitive domain is
\[
 \mathcal H_{\rm prim}
   :=a_K^{-1}(B(K)_{\mathbf Q}).
                                                               \tag{H9.52}
\]
Every class in this domain is admissible for the geometric height.
The normalized signed middle height
$h_m=(-1)^m H_{A_K}$ is positive semidefinite on
$\mathcal H_{\rm prim}$, with exact radical
\[
 \operatorname{rad}(h_m|_{\mathcal H_{\rm prim}})
      =a_K^{-1}(B(k)_{\mathbf Q}).                           \tag{H9.53}
\]
It induces a positive definite form on the quotient
\[
 \frac{\mathcal H_{\rm prim}}
       {a_K^{-1}(B(k)_{\mathbf Q})}
      \ \simeq\
 \frac{B(K)_{\mathbf Q}}{B(k)_{\mathbf Q}}
      \ \simeq\
 \operatorname{Hom}_k(J(C),B)\otimes\mathbf Q.               \tag{H9.54}
\]
The same assertion holds on the image of an actual primitive middle
Chow projector for $A$, with the radical obtained by intersection
with (H9.53). A normalized projected universal family makes its
quotient in (H9.54) surjective as well.
\end{enumerate}
\end{theorem}

\begin{proof}
JAC-P15 proves rational equivalence equals numerical equivalence on
$X_0$ with rational coefficients. Thus the numerical variant of the
spreading lemma applies in every codimension and proves part (1).
For an integral class, membership in the rational span of the
algebraically trivial subgroup means that a nonzero integer multiple
belongs to that subgroup; clear denominators and any remaining
torsion. The converse follows from numerical intersection numbers
having no integer torsion. This proves the $\tau$ formulation.

If $f:(A_0)_F\to A'$ is an isogeny and $z$ is numerically trivial
on $A'$, the projection formula shows that $f^*z$ is numerically
trivial and hence algebraically trivial by the result just proved.
Proper pushforward preserves algebraic
equivalence, and $f_*f^*z=\deg(f)z$. Dividing by $\deg(f)>0$
proves the isogeny variant, including inseparable isogenies.
If the isogeny is only given over $\overline F$, it is defined
over a finite extension $L/F$. Apply the preceding argument over
$L$ to $z_L$, which remains numerically trivial by the projection
formula against norms of complementary cycles. The rational
finite-extension descent argument below then shows that $z$ is
algebraically trivial over $F$. This proves the geometric-isogeny
version as stated.

For part (2), take $F=K$ in part (1). Every geometrically
homologically trivial class is algebraically trivial over $K$ with
rational coefficients. Kahn's Theorem 5.6 puts all of these classes
in the saturated admissible height domain. The primitive incidence
polarization supplied by JAC-T4 remains ample over $K$.
H9-T9 therefore gives positive semidefiniteness and identifies the
radical with the inverse image of the $K/k$-trace of $B_K$.

Since $B_K$ is constant and $k$ is algebraically closed, its trace is
$B$, with the identity trace morphism. This assertion concerns the
trace \emph{variety}; its group of trace points in $B(K)$ is
$B(k)$, which need not equal $B(K)$. Thus H9-T9 gives exactly
(H9.53), not the assertion that every height is zero.

For completeness the first arrow of (H9.54) is surjective.
A normalized rational universal incidence family on $B\times A$
induces the identity of $B$ in the isogeny category. For $b\in B(K)$
its pointed fibre difference is an algebraically trivial cycle with
incidence image $b$, after clearing the single normalization
denominator. Applying the primitive projector to this family gives
the same surjectivity on the primitive Chow image. This is the same
normalized-family construction used in H9-T9, and no surjectivity
from arbitrary homological cycles is being assumed independently.

To prove the last isomorphism of (H9.54), choose $c_0\in C(k)$.
A point $b\in B(K)$ extends uniquely to a morphism $f_b:C\to B$:
properness extends it over every discrete valuation ring on the
smooth curve. The map $f_b-f_b(c_0)$ factors uniquely through the
Abel map $C\to J(C)$. Hence
\[
 B(K)=B(k)\oplus\operatorname{Hom}_k(J(C),B)
\]
as abelian groups after the chosen base point, and the quotient
isomorphism is independent of the constant translation. Tensoring
with $\mathbf Q$ gives (H9.54). In particular a nonzero homomorphism
$J(C)\to B$ is a positive direction in this quotient.

One can also see its height normalization directly. For an integral
point $b$ represented by $f_b:C\to B$, pull the normalized universal
family back to $C\times A$. All fibres are homologically trivial,
so its closure on this smooth product model is admissible. The
definition of the incidence divisor gives
\[
 h_m(z_b,z_b)
      =(-1)^m\deg_C f_b^*\Delta^*d_B
      =2\deg_C f_b^*\ell_B.                               \tag{H9.55}
\]
Rigidification removes constant-family terms; equivalently their
pairings vanish by the incidence-kernel/trace argument in H9-T9.
For rational points, clear denominators first and divide the two
sides by their square. Formula (H9.55) is positive for a nonconstant
$f_b$ and zero for a constant map, because $\ell_B$ is ample.
It is consistent with (H9.53)--(H9.54).
\end{proof}

\paragraph{Why passage through an algebraic closure causes no hidden descent assumption.}
The proof of the spreading lemma already constructs its two parameter
points over $F$ itself. It therefore bypasses any assumption that
integral algebraic equivalence descends from $\overline F$.
Here is the alternative rational norm argument, to specify its scope.
An algebraic-equivalence witness over $\overline F$ is defined over
a finite extension $L/F$. Achter--Casalaina-Martin--Vial, Lemma 4.14,
author pp.~14--15, applied with rational coefficients, replaces that
witness by a family $Z$ on a smooth geometrically integral curve
$T/F$, with two $L$-points $t_0,t_1$. This statement includes
inseparable $L/F$ because the characteristic is inverted. Proper
pushforward of the two fibres gives
\[
 [L:F]z
  =Z_{\beta},\qquad
 \beta=(t_1)_*[\operatorname{Spec}L]
            -(t_0)_*[\operatorname{Spec}L],\qquad
 \deg_F\beta=0.
\]
Lemma 3.9 and the proof of Lemma 4.15 in that source show that
$Z_\beta$ is algebraically trivial over $F$. Division by $[L:F]$
proves rational descent. This is precisely a degree-zero cycle
norm argument; it does not assert integral descent over imperfect
fields or treat an inseparable parameter curve as smooth over $F$.

\paragraph{The isotrivial primitive height corollary.}
More generally let $A'/K$, with $K=k(C)$ as above, be an abelian
variety of dimension $2m-1$ that is geometrically isogenous to an
abelian variety over $\overline{\mathbf F}_p$. Choose a $K$-defined
ample class and let $B'/K$ be its primitive higher Picard factor.
Part (1) identifies the entire geometrically homologically trivial
rational Chow group of $A'$ with its algebraically trivial group
over $K$. JAC-T4 supplies the signed primitive incidence
polarization over $\overline K$; this is the $K$-defined rational
incidence class, so its ampleness descends to $K$. H9-T9 now proves
positive semidefiniteness on
\[
 a_K^{-1}\bigl(B'(K)_{\mathbf Q}\bigr)
 \subset\operatorname{CH}^m(A')_{\rm hom,\mathbf Q},
\]
with exact radical
\[
 a_K^{-1}\!\left(
    \operatorname{im}\bigl(\operatorname{Tr}_{K/k}(B')(k)
                  \longrightarrow B'(K)\bigr)_{\mathbf Q}
             \right).
\]
The normalized universal family makes the quotient the corresponding
group of $K$-points modulo trace points. This argument allows a
finite inseparable field of definition for the geometric isogeny.
It does not assert that $B'$ is constant over $K$ or replace its
trace by all points that become constant after a finite cover.
The final $\operatorname{Hom}_k(J(C),B)$ description in (H9.54)
is specific to the constant case proved there.

\paragraph{Other constant varieties and exact remaining limits.}
The same height conclusion applies to a constant odd-dimensional
$X_0/k_0$ of abelian type whenever an algebraic primitive factor and
its signed incidence polarization have been established: use part (1)
to supply the homological-to-algebraic domain equality and H9-T9 to
supply the height and trace statement. No polarization for an
unspecified abelian-type variety is inferred from the phrase
``abelian type'' alone.

The equalities in (H9.49a) do not assert rational equivalence equals
numerical equivalence after transcendental extension. Even for
abelian varieties, the theorem concerns varieties isogenous
to constants from $\overline{\mathbf F}_p$. It does not apply to a
general abelian variety with transcendental moduli, nor does it
annul every algebraically trivial cycle after extension of constants.
For example a constant elliptic curve has nonconstant points over
its own function field. Their degree-zero divisors are homologically
and algebraically trivial, while their signed geometric self-height
is positive by (H9.55). The radical includes the incidence kernel
and the preimage of constant trace points, not all incidence points.

This gives a valid full homological domain for the corrected H9(b)
height/trace statement in this constant finite-field situation. It
does not repair the invalid homological-to-incidence map for arbitrary
pencils, and it has no implication for the whole arithmetic
$\operatorname{CH}(N)$ of the number-field Fermat Jacobian.

\paragraph{Sources and status.}
The new deduction is the combination of JAC-P15 with the displayed
spreading lemma and the already proved primitive height theorem.
It uses the \emph{Tate} consequence032 Corollary 8.3 through
Kahn2003, Theorem 1.10, published p.~982, and uses the separate
Hodge-standard consequence032 Corollary 8.4 for the primitive
incidence polarization. The parameter-space source was directly read
at \url{https://www.math.uni-bielefeld.de/~vial/abtriv.pdf}:
Definition 3.1 and Proposition 3.5 pp.~4--5, Lemma 3.8 p.~6,
Lemma 3.9 pp.~6--7, Theorem 4.11 and Lemmas 4.14--4.15 pp.~14--16,
and Theorem A.1 pp.~17--18. These are author-manuscript page numbers.
No historical originality or formal verification is asserted.
